% ============================================================================
% LANGENTWURF LE-05d: Mi tiempo libre - Wiederholung + Stundenleistung
% Spanisch Klasse 7 - Sequenz 5: Mi tiempo libre
% ============================================================================
% Tier: 1 (vollständig)
% Doppelstunde: 05d (Std. 7-8 von 8) - ABSCHLUSS
% Fachfarbe: #c41e3a (Karminrot)
% ============================================================================

\documentclass[11pt,a4paper]{article}

\usepackage[utf8]{inputenc}
\usepackage[T1]{fontenc}
\usepackage[ngerman]{babel}
\usepackage{geometry}
\usepackage{fancyhdr}
\usepackage{lastpage}
\usepackage{xcolor}
\usepackage{tcolorbox}
\usepackage{enumitem}
\usepackage{tabularx}
\usepackage{longtable}
\usepackage{array}
\usepackage{booktabs}
\usepackage{amssymb}

% ============================================================================
% KONFIGURATION
% ============================================================================

\newcommand{\fach}{Spanisch}
\newcommand{\fachfarbe}{c41e3a}
\newcommand{\jahrgangsstufe}{7}
\newcommand{\schulart}{Gesamtschule}
\newcommand{\stundenthema}{Mi tiempo libre - Wiederholung + Stundenleistung}
\newcommand{\reihenthema}{Mi tiempo libre}
\newcommand{\stundennummer}{05d}
\newcommand{\reihenstunden}{4}
\newcommand{\gession}{A1}
\newcommand{\lehrwerk}{Qué pasa 1}
\newcommand{\lehrwerkseite}{S. 70-87}
\newcommand{\lerngruppengroesse}{24}

% ============================================================================
% LAYOUT
% ============================================================================
\geometry{left=2.5cm, right=2.5cm, top=2.5cm, bottom=2.5cm}
\definecolor{fach}{HTML}{\fachfarbe}
\definecolor{gray}{HTML}{666666}
\definecolor{lightgray}{HTML}{f5f5f5}
\definecolor{successgreen}{HTML}{2a7a4b}
\definecolor{warningorange}{HTML}{b35c00}
\definecolor{basis}{HTML}{2a7a4b}
\definecolor{standard}{HTML}{2c5aa0}
\definecolor{erweiterung}{HTML}{7b1fa2}

\tcbuselibrary{skins,breakable}

\newtcolorbox{infobox}[1][]{
    colback=lightgray,
    colframe=fach,
    fonttitle=\bfseries,
    title=#1,
    breakable
}

\newtcolorbox{scorebox}{
    colback=white,
    colframe=fach,
    fonttitle=\bfseries,
    title=Didaktik-Score,
    breakable
}

\pagestyle{fancy}
\fancyhf{}
\fancyhead[L]{\small\textcolor{gray}{\fach{} -- Jg. \jahrgangsstufe{} -- \stundenthema}}
\fancyhead[R]{\small\textcolor{gray}{LE-\stundennummer{} | Tier 1}}
\fancyfoot[C]{\small\thepage{} / \pageref{LastPage}}
\renewcommand{\headrulewidth}{0.4pt}

% Differenzierungsmarker
\newcommand{\B}{\textcolor{basis}{\textbf{[B]}}}
\newcommand{\Std}{\textcolor{standard}{\textbf{[S]}}}
\newcommand{\E}{\textcolor{erweiterung}{\textbf{[E]}}}

% ============================================================================
% DOKUMENT
% ============================================================================
\begin{document}

% ============================================================================
% DECKBLATT
% ============================================================================
\begin{titlepage}
    \centering
    \vspace*{2cm}
    {\large\textcolor{gray}{\schulart}}\\[2cm]
    {\Huge\textcolor{fach}{\textbf{Langentwurf}}}\\[0.5cm]
    {\large LE-\stundennummer{} | Tier 1 (vollständig)}\\[2cm]
    {\LARGE\textbf{\stundenthema}}\\[0.5cm]
    {\large Sequenz: \reihenthema{} -- \textbf{ABSCHLUSS}}\\[2cm]
    \begin{tabular}{rl}
        \textbf{Fach:} & \fach{} (A1) \\[0.3cm]
        \textbf{Klasse:} & \jahrgangsstufe \\[0.3cm]
        \textbf{Lehrwerk:} & \lehrwerk, \lehrwerkseite \\[0.3cm]
        \textbf{Doppelstunde:} & \stundennummer{} (Std. 7-8 von 8) \\[0.3cm]
    \end{tabular}
    \vfill
\end{titlepage}

\tableofcontents
\newpage

% ============================================================================
% 1. BEDINGUNGSANALYSE
% ============================================================================
\section{Bedingungsanalyse}

\subsection{Lerngruppe}
\begin{tabular}{ll}
    Kursgröße: & \lerngruppengroesse{} SuS \\
    Jahrgangsstufe: & \jahrgangsstufe \\
    Sprachniveau: & A1 (GER) -- Anfänger \\
    Fremdsprache: & 2. Fremdsprache (nach Englisch) \\
\end{tabular}

\subsection{Vorwissen aus 05a/05b/05c}
\begin{itemize}
    \item \textbf{05a:} Hobbys und Freizeitaktivitäten (jugar, tocar, hacer, ver)
    \item \textbf{05b:} Verben auf -er und -ir (comer, beber, leer, escribir, vivir)
    \item \textbf{05c:} gustar + Infinitiv (Me gusta..., Te gusta..., Le gusta...)
\end{itemize}

\subsection{Differenzierung (intern)}
\begin{tabular}{|l|p{3.5cm}|p{3.5cm}|p{3.5cm}|}
\hline
& \B{} \textbf{Basis} & \Std{} \textbf{Standard} & \E{} \textbf{Erweiterung} \\
\hline
\textbf{Ziel} & BR: Rezeption, Grundwortschatz & MR: Produktion mit Hilfe & GY: Freie Anwendung \\
\hline
\textbf{Verben} & Konjugation ablesen & Konjugation bilden & Freie Satzbildung \\
\hline
\textbf{Test} & Basisaufgaben (60\%) & Alle Pflichtaufgaben & + Zusatzaufgaben \\
\hline
\end{tabular}

% ============================================================================
% 2. SACHANALYSE
% ============================================================================
\section{Sachanalyse}

\subsection{Sprachliche Strukturen: Verben auf -er/-ir}

\textbf{Konjugation -er (z.B. comer = essen):}
\begin{center}
\begin{tabular}{|l|l|l|}
\hline
\textbf{Person} & \textbf{Endung} & \textbf{comer} \\
\hline
yo & -o & como \\
tú & -es & comes \\
él/ella & -e & come \\
nosotros & -emos & comemos \\
vosotros & -éis & coméis \\
ellos/ellas & -en & comen \\
\hline
\end{tabular}
\end{center}

\textbf{Konjugation -ir (z.B. vivir = leben):}
\begin{center}
\begin{tabular}{|l|l|l|}
\hline
\textbf{Person} & \textbf{Endung} & \textbf{vivir} \\
\hline
yo & -o & vivo \\
tú & -es & vives \\
él/ella & -e & vive \\
nosotros & -imos & vivimos \\
vosotros & -ís & vivís \\
ellos/ellas & -en & viven \\
\hline
\end{tabular}
\end{center}

\subsection{gustar + Infinitiv}
\begin{center}
\begin{tabular}{|l|l|l|}
\hline
\textbf{Person} & \textbf{Indirektes Pronomen} & \textbf{Beispiel} \\
\hline
ich & me & Me gusta jugar al fútbol. \\
du & te & ¿Te gusta leer? \\
er/sie & le & Le gusta escuchar música. \\
wir & nos & Nos gusta bailar. \\
ihr & os & ¿Os gusta cantar? \\
sie & les & Les gusta ver la tele. \\
\hline
\end{tabular}
\end{center}

\subsection{Wichtige Hobbys (Wortschatz)}
\begin{itemize}
    \item \textbf{jugar:} jugar al fútbol, jugar al baloncesto, jugar a videojuegos
    \item \textbf{tocar:} tocar la guitarra, tocar el piano
    \item \textbf{hacer:} hacer deporte, hacer los deberes
    \item \textbf{escuchar:} escuchar música
    \item \textbf{ver:} ver la tele, ver películas
    \item \textbf{leer:} leer libros, leer cómics
\end{itemize}

% ============================================================================
% 3. DIDAKTISCHE ANALYSE
% ============================================================================
\section{Didaktische Analyse}

\subsection{Stellung in der Sequenz}
Dies ist die \textbf{4. und letzte Doppelstunde} (Std. 7-8 von 8) der Sequenz ``\reihenthema''.
Die Stunde schließt die Sequenz ab mit:
\begin{enumerate}
    \item Konsolidierung aller bisherigen Inhalte (Hobbys, -er/-ir Verben, gustar)
    \item Interaktiver Wiederholung (Lernpfad)
    \item Stundenleistung (summative Bewertung)
\end{enumerate}

\subsection{Kompetenzziele (Rahmenplan MV)}
\begin{itemize}
    \item \textbf{Sprechen:} Über Hobbys und Vorlieben sprechen
    \item \textbf{Schreiben:} Kurze Texte über Freizeitaktivitäten verfassen
    \item \textbf{Verfügen über sprachliche Mittel:} -er/-ir-Konjugation, gustar-Konstruktion
\end{itemize}

\subsection{Legitimation}
\textbf{Rahmenplan Spanisch MV:}
\begin{itemize}
    \item An Gesprächen teilnehmen: über Freizeitaktivitäten sprechen
    \item Verfügen über sprachliche Mittel: regelmäßige Verben konjugieren
\end{itemize}

% ============================================================================
% 4. STUNDENZIELE
% ============================================================================
\section{Stundenziele}

\subsection{Grobziel}
Die SuS \textbf{konsolidieren} ihre Kenntnisse über Hobbys, -er/-ir-Verben und gustar und weisen ihre Kompetenzen in einer \textbf{Stundenleistung} nach.

\subsection{Feinziele (operationalisiert)}
\begin{tabular}{|c|p{8cm}|c|c|}
\hline
\textbf{Nr.} & \textbf{Die SuS können...} & \textbf{AFB} & \textbf{Diff.} \\
\hline
FZ1 & Hobbys benennen und kategorisieren (jugar, tocar, hacer, ver) & I & alle \\
\hline
FZ2 & regelmäßige Verben auf -er und -ir korrekt konjugieren & I/II & alle \\
\hline
FZ3 & mit gustar + Infinitiv Vorlieben ausdrücken & II & alle \\
\hline
FZ4 & einen kurzen Text über eigene Hobbys verfassen & II/III & \Std{}/\E{} \\
\hline
\end{tabular}

% ============================================================================
% 5. METHODISCHE ANALYSE
% ============================================================================
\section{Methodische Analyse}

\subsection{Medien}
\begin{itemize}
    \item Tafel/Whiteboard
    \item Lehrwerk: \lehrwerk, \lehrwerkseite
    \item Arbeitsblatt: AB-05d.pdf (Wiederholung)
    \item Stundenleistung: SL-05d.pdf (Test)
    \item Lernpfad: LP-05d.html (digital)
    \item Tablets/Computer für LP (optional)
\end{itemize}

\subsection{Sozialformen}
\begin{itemize}
    \item Plenum: Wiederholung, Warm-up
    \item Partnerarbeit: Dialogübungen
    \item Einzelarbeit: Stundenleistung
\end{itemize}

\subsection{Besonderheit: Stundenleistung}
\begin{itemize}
    \item Dauer: 25-30 Minuten
    \item Umfang: 20 Punkte (4 Aufgaben)
    \item Inhalte: Hobbys, -er/-ir Verben, gustar
    \item Bewertung: Note nach Prozentskala
\end{itemize}

% ============================================================================
% 6. VERLAUFSPLANUNG
% ============================================================================
\section{Verlaufsplanung}

\textbf{1. Stunde (45 Min.) -- Schwerpunkt: Wiederholung}

\begin{longtable}{|p{1cm}|p{1.5cm}|p{4cm}|p{3cm}|p{2cm}|p{2cm}|}
\hline
\textbf{Zeit} & \textbf{Phase} & \textbf{Lehrerhandeln} & \textbf{Schülerhandeln} & \textbf{SF} & \textbf{Medien} \\
\hline
\endhead

0--5 & Einstieg & ``¿Qué te gusta hacer en tu tiempo libre?'' & Antworten sammeln & Plenum & Tafel \\
\hline

5--15 & Wiederholung & Hobbys kategorisieren (jugar, tocar, hacer) & Zuordnen, Ergänzen & Plenum & Tafel \\
\hline

15--25 & Übung 1 & AB-05d Aufg. 1-2 (-er/-ir Konjugation) & Bearbeiten & EA & AB \\
\hline

25--35 & Übung 2 & AB-05d Aufg. 3 (gustar) & Sätze bilden & EA/PA & AB \\
\hline

35--45 & Sicherung & Besprechung, Merkkasten & Korrigieren, Notieren & Plenum & AB \\
\hline
\end{longtable}

\vspace{1em}

\textbf{2. Stunde (45 Min.) -- Schwerpunkt: Stundenleistung}

\begin{longtable}{|p{1cm}|p{1.5cm}|p{4cm}|p{3cm}|p{2cm}|p{2cm}|}
\hline
\textbf{Zeit} & \textbf{Phase} & \textbf{Lehrerhandeln} & \textbf{Schülerhandeln} & \textbf{SF} & \textbf{Medien} \\
\hline
\endhead

0--5 & Aktivierung & Schnelle Wiederholung: 3 Verben konjugieren & Antworten & Plenum & -- \\
\hline

5--10 & Vorbereitung & Stundenleistung erklären, Fragen klären & Zuhören, Fragen & Plenum & SL \\
\hline

10--35 & Test & Aufsicht, individuelle Hilfe bei Verständnisfragen & Stundenleistung bearbeiten & EA & SL \\
\hline

35--40 & Abgabe & Einsammeln, kurze Entspannung & Abgeben & -- & -- \\
\hline

40--45 & Abschluss & Ausblick Sequenz 6, Feedback & Rückmeldung & Plenum & -- \\
\hline
\end{longtable}

\textbf{Legende:} EA = Einzelarbeit, PA = Partnerarbeit, SF = Sozialform, SL = Stundenleistung

% ============================================================================
% 7. ERWARTETE SCHWIERIGKEITEN
% ============================================================================
\section{Erwartete Schwierigkeiten}

\begin{tabular}{|p{5cm}|p{8cm}|}
\hline
\textbf{Schwierigkeit} & \textbf{Reaktion} \\
\hline
Verwechslung -er/-ir Endungen & Unterschied betonen: -emos vs. -imos \\
\hline
gustar-Konstruktion (Subjekt/Objekt) & Vereinfachte Erklärung: ``mir gefällt'' \\
\hline
Prüfungsangst bei Stundenleistung & Ruhige Atmosphäre, klare Instruktionen \\
\hline
Wortschatz Hobbys vergessen & Wortschatzliste an der Tafel \\
\hline
\end{tabular}

% ============================================================================
% 8. HAUSAUFGABE
% ============================================================================
\section{Hausaufgabe}

\begin{itemize}
    \item \textbf{Vor der Stunde:} LP-05d.html zur Wiederholung (optional)
    \item \textbf{Nach der Stunde:} Vokabeln Sequenz 5 wiederholen
\end{itemize}

% ============================================================================
% 9. LÖSUNGEN
% ============================================================================
\section{Lösungen AB-05d}

\textbf{Aufgabe 1 (-er Verben):} (4P)
\begin{itemize}
    \item yo como, tú bebes, él/ella lee, nosotros comemos
\end{itemize}

\textbf{Aufgabe 2 (-ir Verben):} (4P)
\begin{itemize}
    \item yo vivo, tú escribes, él/ella vive, nosotros vivimos
\end{itemize}

\textbf{Aufgabe 3 (gustar):} (6P)
\begin{itemize}
    \item Me gusta jugar al fútbol.
    \item Te gusta escuchar música.
    \item Le gusta leer libros.
\end{itemize}

\textbf{Aufgabe 4 (Hobbys zuordnen):} (6P)
\begin{itemize}
    \item jugar: al fútbol, al baloncesto, a videojuegos
    \item tocar: la guitarra, el piano
    \item hacer: deporte, los deberes
\end{itemize}

% ============================================================================
% 10. STUNDENLEISTUNG - ÜBERSICHT
% ============================================================================
\section{Stundenleistung SL-05d}

\textbf{Aufbau:} 4 Aufgaben, 20 Punkte, 25 Minuten

\begin{tabular}{|c|l|c|c|c|}
\hline
\textbf{Aufg.} & \textbf{Inhalt} & \textbf{AFB} & \textbf{Punkte} & \textbf{Diff.} \\
\hline
1 & Hobbys zuordnen (jugar/tocar/hacer) & I & 4 & alle \\
\hline
2 & -er/-ir Verben konjugieren & I/II & 6 & alle \\
\hline
3 & gustar-Sätze vervollständigen & II & 4 & alle \\
\hline
4 & Kurzer Text über Hobbys & II/III & 6 & alle \\
\hline
\multicolumn{3}{|r|}{\textbf{Gesamt:}} & \textbf{20} & \\
\hline
\end{tabular}

\textbf{Bewertung (Klasse 7):}
\begin{tabular}{|c|c|c|c|c|c|c|}
\hline
Note & 1 & 2 & 3 & 4 & 5 & 6 \\
\hline
ab \% & 96 & 80 & 60 & 40 & 20 & <20 \\
\hline
ab Punkte & 19 & 16 & 12 & 8 & 4 & <4 \\
\hline
\end{tabular}

% ============================================================================
% 11. DIDAKTIK-SCORE
% ============================================================================
\newpage
\section{Didaktik-Score}

\begin{scorebox}
\textbf{Bewertung der didaktischen Qualität nach 10 Dimensionen}\\
\textbf{Ziel-Score: $\geq$ 9.0 (Premium-Qualität)}

\vspace{1em}
\begin{tabular}{|c|l|c|c|p{4.5cm}|}
\hline
\textbf{\#} & \textbf{Dimension} & \textbf{Gew.} & \textbf{Score} & \textbf{Notiz} \\
\hline
1 & Differenzierung & 15\% & 9/10 & Basis-, Standard-, Erweiterungs-Niveau \\
\hline
2 & Taxonomie (AFB) & 10\% & 9/10 & AFB I: 35\%, AFB II: 45\%, AFB III: 20\% \\
\hline
3 & Lernziel-Alignment & 10\% & 10/10 & Rahmenplan: Freizeitaktivitäten \\
\hline
4 & Aktivierung & 10\% & 9/10 & Wiederholung, PA, Test \\
\hline
5 & Scaffolding & 10\% & 9/10 & Übung $\rightarrow$ Anwendung $\rightarrow$ Test \\
\hline
6 & Feedback-Qualität & 10\% & 9/10 & Besprechung vor Test \\
\hline
7 & Sprachliche Klarheit & 10\% & 10/10 & Altersgerecht Kl. 7 \\
\hline
8 & Motivation \& Relevanz & 10\% & 10/10 & Hobbys = hohe Alltagsrelevanz \\
\hline
9 & Inklusion (UDL) & 10\% & 9/10 & Konjugationstabelle als Hilfe \\
\hline
10 & Praktikabilität & 5\% & 9/10 & 45+45 Min realistisch \\
\hline
\multicolumn{3}{|r|}{\textbf{GESAMT:}} & \textbf{9.2/10} & \textcolor{successgreen}{\textbf{FREIGABE}} \\
\hline
\end{tabular}
\end{scorebox}

\end{document}
